\expandafter\def\csname HistoricalNoteText550\endcsname{Zuweilen (obwohl äusserst selten) sind uns allerdings Fälle vorgekommen, die eine Ausnahme davon machten, und wo Spuren von Beschleunigung oder Retardation der Änderung in so kurzen Zwischenzeiten sich doch unverkennbar nachweisen liessen. Mit Ausführlichkeit soll dieser Gegenstand in Zukunft abgehandelt werden,}
\expandafter\def\csname HistoricalNoteText544\endcsname{\enspace\resizebox{.925\linewidth}{!}{Im 9. Bande der Astronomischen Nachrichten hat dieser hochverdiente Naturforscher uns auch mit}}
\par\clearpage
\SourceLeaf{542}\GaussVSourceFolio{536}{L}{MAGNETISCHE BEOBACHTUNGEN.}
\noindent als die Kette offen ist, fängt gleichfalls in dem Augenblick an mitzuschwingen, wo die Kette geschlossen wird, allein diese Schwingungen, von derselben Dauer, wie die natürlichen, nehmen an Grösse beständig zu, bis diese (erst nach sehr langer Zeit) zu einem Maximum kommt, wo der Widerstand der Luft der Vergrösserung durch die Inductionskraft das Gleichgewicht hält. Um diesen merkwürdigen Versuch wirklich anstellen zu können, wurde (da die Aufhängung eines grossen Stabes wegen Mangel eines zweiten dafür passenden Multiplicators jetzt nicht thunlich war) der einpfündige Stab des physicalischen Cabinets durch Verbindung mit einem ähnlichen etwas schwächer magnetisirten auf bekannte Weise astatisch gemacht, oder vielmehr zu einer Doppelnadel, deren natürliche Schwingungsdauer genau auf $42''3$ gebracht wurde. Der Versuch gelang damit auf das vollkommenste. Der in der Sternwarte schwingende Stab theilte dieser Doppelnadel im physicalischen Cabinet, in dem Augenblick wo die Kette geschlossen wurde, wie durch eine wunderbare Sympathie seine Schwingungen mit, und zwar so, dass jede folgende etwa 50 Scalentheile oder einen halben Grad grösser wurde, als die vorhergehende. Bald ging das ganze Scalenbild aus dem Felde, allein fortwährend konnte man an der immer wachsenden Schnelligkeit, mit welcher das Scalenbild durch das Gesichtsfeld ging, die Zunahme des Schwingungsbogens erkennen. Über eine Stunde wurde dies wunderbar sympathetische Spiel beobachtet.

Es braucht kaum bemerkt zu werden, dass auch der vierpfündige Stab im magnetischen Observatorium in die geschlossene Kette einen Strom inducirt, dessen Dasein an der schnellen Abnahme des Schwingungsbogens auf das bestimmteste erkannt wird, und der daher auch auf die beiden andern Stäbe Wirkungen ausüben muss, denen ähnlich, welche der erstere Versuch gezeigt hat; allein die Rechnung ergibt, und die Erfahrung bestätigt, dass diese Wirkungen zu klein ausfallen, um merklich zu sein. Noch weniger könnte also der schwächste Stab unter den dreien merkliche Wirkungen dieser Art erzeugen.

\par\vspace{.85\baselineskip}
\centerline{\rule{.255\linewidth}{.35pt}}
\clearpage
\SourceLeaf{543}\GaussVSourceFolio{537}{R}{MAGNETISCHE BEOBACHTUNGEN}
\pdfbookmark[0]{Magnetnadel in Copenhagen und Mailand. 1834}{copenhagen_mailand_1834}
\begin{center}
\resizebox{.939\linewidth}{!}{{\itshape Beobachtungenen der Variationen der Magnetnadel in Copenhagen und Mailand}}\par
\resizebox{.356\linewidth}{!}{{\itshape am \textup{5.} und \textup{6.} November \textup{1834.}}}\par
\vspace{.19in}
\begin{minipage}{.779\linewidth}
\centering
\hrule height .35pt
\vspace{1.55mm}
\noindent\resizebox{.810\linewidth}{!}{\small \textsc{Schumacher}.\quad Astronomische Nachrichten Nr.\ 276.\quad 1835 März 21.}\par
\vspace{1.55mm}
\hrule height .35pt
\end{minipage}
\end{center}

Seit der Vollendung des hiesigen magnetischen Observatoriums werden hier unter andern regelmässig an gewissen im Voraus bestimmten Tagen die Variationen der magnetischen Declination ununterbrochen in kurzen Zeitintervallen beobachtet, wozu anfangs dieselben Termine gewählt waren, welche Herr \textsc{von Humboldt} schon vor mehreren Jahren angeordnet hatte. Seit dem vorigen Frühjahr haben sich schon ziemlich viele Astronomen und Physiker in den Besitz von ähnlichen Apparaten gesetzt, wie der hiesige ist, den ich an einem andern Orte hinlänglich beschrieben habe, und nehmen an jenen verabredeten Beobachtungen Theil. Gleich die ersten auf diese Weise gewonnenen gleichzeitigen Beobachtungen am 4. und 5. Mai, in Göttingen und Waltershausen (einem Gute in der Gegend von Schweinfurt, wo Herr \textsc{Sartorius} mit einem zwar kleinen, aber sonst dem hiesigen ganz ähnlichen Apparat beobachtete), zeigten eine überaus merkwürdige Harmonie in dem vielfach hin und her springenden Gange der Variationen, nicht blos in den grössern sondern auch in den geringern. Ähnliche Erfolge haben sich seitdem in den spätern Terminen, wo Leipzig, Berlin, Braunschweig und Copenhagen Theil genommen haben, schon vielfach wiederholt; einige Proben sind in graphischen Darstellungen in \textsc{Poggendorffs} Annalen mitgetheilt.

So wie sich die Theilnahme an diesen verabredeten Beobachtungen immer weiter verbreiten wird, stehen natürlich immer interessantere und fruchtbarere Resultate zu erwarten. Ich wiederhole daher hier die bereits anderwärts gemachte Anzeige, dass wir, seitdem die Nothwendigkeit, in \textit{sehr} kurzen Zeitintervallen zu beobachten, sich so klar herausgestellt hat, uns veranlasst gefunden haben, mit den Terminen eine Abänderung zu treffen, indem wir die Anzahl der Termine von 8 auf 6 im Jahr, und ihre Dauer von 44 auf 24 Stunden herabgesetzt haben. Die gegenwärtige Bestimmung ist der letzte Sonnabend jedes ungeraden Monats, vom Göttinger Mittag an bis zum Mittag des folgenden Tages. Es kommen zu diesen Hauptterminen noch jedesmal zwei Nebentermine, nemlich am
\GaussVSourceGathering{68}
\par\clearpage
\begingroup
\renewcommand{\footnoterule}{\kern-4.35pt\hbox{\hspace{.012\columnwidth}\rule{.255\columnwidth}{.35pt}}\kern4pt}
\SourceLeaf{544}\GaussVSourceFolio{538}{L}{MAGNETISCHE BEOBACHTUNGEN.}
\noindent nächstfolgenden Dinstag und Mittwoch Abends von 8 bis 10 Uhr. Umständlichere Nachricht, auch über Beobachtungsweise, findet man in \textsc{Poggendorffs} Annalen Bd.\ 33. [S.\ 528 d.\ B.]

Das merkwürdigste bisher erhaltene Resultat bieten die gleichzeitigen Beobachtungen von Copenhagen und Mailand dar, am 5. und 6. November d.\ J., einem Termine nach dem frühern Arrangement, von dessen Abänderung die Beobachter an jenen Orten die Nachricht noch nicht erhalten hatten. In Copenhagen, wo jetzt unter Leitung des Herrn Etatsrath \textsc{Oersted} ein dem hiesigen ganz ähnliches magnetisches Observatorium errichtet ist, wurde eine Nadel von derselben Stärke, wie die hiesige, gebraucht (vier Pfund schwer); in Mailand beobachteten auf der dortigen Sternwarte die Herrn \textsc{Sartorius} und Doctor \textsc{Listing}, unter Beistand des Herrn \textsc{Kreil}, Eleven der Sternwarte, mit der schon oben erwähnten kleinern Nadel. Ich gestehe, dass ich, auch nach den vielen schon früher vorgekommenen Erfahrungen ähnlicher Art, doch durch die Grösse der Übereinstimmung an zwei mehr als 150 Meilen von einander entfernten Orten überrascht wurde. Der blosse Anblick der beigefügten graphischen Darstellung spricht hier für sich. Ich begleite dieselbe nur mit einigen Erläuterungen und Bemerkungen.

Da mir anfangs der Werth der Scalentheile in Copenhagen noch unbekannt war, so entwarf ich die Zeichnung nach solchem Maassstabe, dass die Anomalien ungefähr gleich gross erscheinen, was ich erhielt, indem ich der Seite der Netzquadrate neun Scalentheile der Copenhagner, und drei der Mailänder Beobachtungen entsprechen liess. Ein Scalentheil in Copenhagen beträgt übrigens $21''576$, einer in Mailand $29''341$. In Bogentheilen waren also die Copenhagner Bewegungen etwa $2{,}2$ mal grösser als die Mailänder. Will man hieraus auf das Verhältniss der dabei thätigen Kräfte schliessen, so muss man nicht übersehen, dass diese Erscheinungen nur als Störungen der horizontalen erdmagnetischen Kraft an beiden Orten zu betrachten sind, und dass an einem Orte, wo letztere kleiner ist, eine gleiche störende Kraft grössere Änderungen hervorbringen muss, als an einem andern, wo jene grösser ist. Das Verhältniss der horizontalen Kraft des Erdmagnetismus in Copenhagen und Mailand schätzte ich nach \textsc{Hansteens} schöner Karte im 7. Bande der Astronomischen Nachrichten\HistoricalNote{544} wie
\par\vfill
\clearpage
\endgroup
\SourceLeaf{545}\GaussVSourceFolio{539}{R}{MAGNETISCHE BEOBACHTUNGEN.}
\noindent 1 zu $1{,}23$; danach würde sich also das Verhältniss der störenden Kräfte, die die beträchtlichsten Anomalien an jenen Tagen in Copenhagen und Mailand hervorgebracht haben, etwa wie $1{,}8$ zu 1 schätzen lassen. Wie viel besser werden wir aber in Zukunft über solche räthselhafte Naturkräfte urtheilen können, wenn erst ähnliche gleichzeitige Beobachtungen an vielen weit von einander entlegenen Orten uns zu Gebote stehen werden.

Neben der überraschend grossen Übereinstimmung in dem Gange der Anomalien bemerken wir allerdings auch Verschiedenheiten. Aber es scheint, dass wir nicht über diese uns zu verwundern haben, sondern vielmehr darüber, dass die Unterschiede vergleichungsweise so klein sind. Wir kennen freilich die Ursachen der Erscheinungen noch gar nicht; aber gerade bei dem bunten Spiel ihres Wechsels scheint es unnatürlich, anzunehmen, dass sie alle von Einem Punkt her wirkten: einige Ursachen mögen hier, andere dort ihren Sitz gehabt haben, und so mögen in den 44 Stunden auch wol manche Kräfte von ganz andern Gegenden her, die ein ganz anderes Verhältniss für die beiden Örter hatten, ihr Spiel eingemischt haben. Dass im Allgemeinen die Curve für Mailand viel krauser erscheint, als die für Copenhagen, erklärt sich übrigens von selbst durch den Umstand, dass an ersterm Orte alle 5 Minuten, an letzterm alle 10 Minuten beobachtet wurde; bei den längern Zwischenzeiten mussten folglich manche kleinere und schneller wechselnde Anomalien unbemerkt bleiben.

Wenngleich das Interesse für diese Forschungen einer Verstärkung nicht bedarf, so glaube ich doch noch einen Umstand hervorheben zu müssen, der die Astronomen noch besonders berührt. Ob die bei diesen Bewegungen thätigen Kräfte eine messbare Zeit gebrauchen, um sich durch grosse Räume fortzupflanzen, wissen wir noch nicht; diese interessante Frage wird aber ohne Zweifel in
\par\vfill
\begingroup
\footnotesize
\noindent\rule{.252\linewidth}{.35pt}\par
\nobreak\vspace{2mm}
\noindent einer allgemeinen Karte für die \textit{ganze} Intensität beschenkt. So dankbar man diese schöne Arbeit anerkennen muss, so kann ich doch die Bemerkung nicht unterdrücken, dass eine allgemeine Karte für die \textit{horizontale} Intensität in \textit{vielfacher} Hinsicht noch ungleich nützlicher sein würde, namentlich auch in Verbindung mit einer zuverlässigen allgemeinen Declinations-Karte, zu einer durchgreifenden Begründung einer allgemeinen Theorie. Zu \textit{diesem} Zweck ist die Bestimmung der magnetischen Kraft durch Angabe der ganzen Intensität, Inclination und Declination (die man wol als die einfachste Wahl der Elemente zu betrachten gewohnt ist) gerade die am wenigsten brauchbare. Die weitere Entwickelung dieser Behauptung, die vielleicht manchem paradox scheinen könnte, muss ich mir aber für einen andern Ort vorbehalten. Möchte nur jener Naturforscher uns aus der Fülle seiner gesammelten Schätze bald mit jenen Erfordernissen beschenken.\par
\endgroup
\GaussVSourceGathering{68\textsuperscript{*}}
\clearpage
\SourceLeaf{546}\GaussVSourceFolio{540}{L}{MAGNETISCHE BEOBACHTUNGEN.}
\noindent Zukunft ihre Beantwortung finden. Ist die Zeit unmessbar klein, so werden solche Beobachtungen schneller auf- und abgehender Bewegungen zu Längenbestimmungen dienen können, die unter vortheilhaften Umständen selbst den schärfern zur Seite gestellt werden dürfen. Aus vorgekommenen Bewegungen in Göttingen und Leipzig habe ich schon mehreremale unter jener Voraussetzung den Längenunterschied auf eine halbe Zeitminute richtig ableiten können. Allein zuweilen zeigen sich so schnelle Bewegungen, dass daraus noch viel schärfere Zeitbestimmungen abgeleitet werden können. Die stärksten Bewegungen, die mir bisher vorgekommen sind, fanden statt am 7. Februar d.\ J., wo den ganzen Tag die Nadel überaus unruhig war. Ich beobachtete Bewegungen von 17 Scalentheilen oder 6 Bogenminuten in Einer Zeitminute, einige Minuten regelmässig andauernd, dann nach und nach langsamer werdend, und nachher in die entgegengesetzte übergehend. Dergleichen Erscheinungen an zwei Orten mit guten Apparaten (die selbst in einzelnen Beobachtungen eine Genauigkeit bis auf wenige Bogensecunden geben) sorgfältig verfolgt, könnten, wenn die Wirkung der Kräfte in unmessbar kleiner Zeit geschieht, den Längenunterschied auf eine Zeitsecunde genau geben. Jedenfalls erhellt, wie wichtig es zur Aufklärung des Gegenstandes sein wird, dass alle Beobachter, denen die Mittel dazu zu Gebote stehen, immer für eine gute Zeitbestimmung Sorge tragen.

Schliesslich bemerke ich noch, dass die Beobachter in Mailand die dortige Inclination mit einem \textsc{Lenoir}schen Inclinatorium am 2. November $=63^0\,55'\,26''$ gefunden haben.

\par\vspace{2.9\baselineskip}
\pdfbookmark[1]{Brief von Gauss an Schumacher. 1836}{letter546}
\begin{center}
\resizebox{.386\linewidth}{!}{{\itshape Brief von Gauss an Schumacher.}}\par
\vspace{.17in}
\begin{minipage}{.724\linewidth}
\centering
\hrule height .35pt
\vspace{1.55mm}
\noindent\resizebox{.887\linewidth}{!}{\small \textsc{Schumacher}.\quad Astronomische Nachrichten Nr.\ 310.\quad 1836. Juni 11.}\par
\vspace{1.55mm}
\hrule height .35pt
\end{minipage}
\end{center}
\noindent\makebox[\linewidth][r]{\small Göttingen. 1836. \quad April 23.\hspace{.064\linewidth}}\par
Es waren heute Morgen ausserordentliche Bewegungen der Magnetnadel, noch grösser als am 7. Februar 1835. Dies veranlasste mich einige Sets in der Sternwarte zu beobachten, während Dr.\ \textsc{Goldschmidt} im magnetischen~Observa-
\par\clearpage
\SourceLeaf{547}\GaussVSourceFolio{541}{R}{MAGNETISCHE BEOBACHTUNGEN.}
\noindent torium aufzeichnete. Der gleichförmige Gang bestätigte sich hier so schön, dass ich es wagte den gegenseitigen Uhrstand daraus abzuleiten.

\par
\begingroup
\newcommand{\GaussVClockReading}[2]{\makebox[.110\linewidth][l]{%
  \hspace{.005\linewidth}\makebox[.030\linewidth][l]{\ensuremath{#1}}\ensuremath{#2}}}
\newcommand{\GaussVClockLeaders}{\makebox[.363\linewidth][l]{%
  .\hfill.\hfill.\hfill.\hfill.\hfill.\hfill.\hfill.\hfill.\hfill.\hspace{.024\linewidth}}}
\noindent\hspace{.073\linewidth}%
\begin{tabular}{@{}l@{}l@{}}
\makebox[.813\linewidth][l]{\resizebox{.802\linewidth}{!}{Es fand sich, aus einem schnellen Aufsteigen Campa vor Shelton}}&
  \GaussVClockReading{4'}{41''1}\\
\makebox[.813\linewidth][l]{\resizebox{.742\linewidth}{!}{Aus einem wenige Minuten nachher erfolgten Niedersteigen}}&
  \GaussVClockReading{4}{42.4}\\[1.35pt]
\cline{2-2}\noalign{\vskip 1.9pt}
\makebox[.813\linewidth][r]{Mittel\hspace{.023\linewidth}}&
  \GaussVClockReading{4'}{41''7}
\end{tabular}\par
\noindent\hspace{.073\linewidth}Eine directe Vergleichung der Uhren gab,\par
\noindent\hspace{.073\linewidth}%
\begin{tabular}{@{}l@{}l@{}l@{}}
\makebox[.450\linewidth][l]{1) durch ein Zeichen am Fenster}&
  \GaussVClockLeaders&\GaussVClockReading{4'}{41''5}\\
\makebox[.450\linewidth][l]{2) durch einen Inductionsimpuls}&
  \GaussVClockLeaders&\GaussVClockReading{4}{41.5}
\end{tabular}\par
\endgroup
\noindent Also eine herrliche Bestätigung dessen, was ich Astronomische Nachrichten Nr.\ 276 [S.\ 539 d.\ B.] gesagt habe.

\par\vspace{3.5\baselineskip}
\pdfbookmark[1]{Das in den Beobachtungsterminen anzuwendende Verfahren.}{verfahren547}
\begin{center}
\resizebox{.708\linewidth}{!}{{\itshape Das in den Beobachtungsterminen anzuwendende Verfahren.}}\par
\vspace{.17in}
\begin{minipage}{.749\linewidth}
\centering
\hrule height .35pt
\vspace{1.55mm}
\noindent\resizebox{.892\linewidth}{!}{\small Resultate aus den Beobachtungen des magnetischen Vereins.\quad 1836.\quad II.}\par
\vspace{1.55mm}
\hrule height .35pt
\end{minipage}
\end{center}

Die sechs jährlich festgesetzten Termine fallen gegen das Ende der Monate Januar, März, Mai, Julius, September, November; sie fangen an am letzten Sonnabend in jedem dieser Monate, Mittags nach Göttinger mittlerer Zeit, und schliessen am Mittag des folgenden Tages; die bisher jedem Haupttermine hinzugefügten Nebentermine (Abends von 8—10 Uhr am Dinstag und Mittwoch der folgenden Woche) werden künftig wegfallen.

In jedem Termine wird, der Regel nach, der Stand der Magnetnadel von fünf zu fünf Minuten bestimmt, so dass ein Termin 289 Resultate gibt. In Göttingen wird die Uhr vor Anfang jedes Termins genau auf mittlere Zeit gestellt. Da eine nahe Gleichzeitigkeit der einzelnen Bestimmungen an den verschiedenen Beobachtungsorten sehr wünschenswerth ist, so haben die Beobachter an den meisten andern Orten die Gewohnheit, ihre Uhren gleichfalls auf Göttinger mittlere Zeit zu stellen. Wo dies nicht wohl geschehen kann, ist zu empfehlen, dass man zu den Beobachtungsmomenten diejenigen vollen Minuten der Uhr wähle, die den Göttinger Beobachtungszeiten am nächsten kommen. Hätte man z.\ B.\ vor
\par\clearpage
\SourceLeaf{548}\GaussVSourceFolio{542}{L}{MAGNETISCHE BEOBACHTUNGEN.}
\noindent Anfang des Termins ausgemittelt, dass die bei der Beobachtung zu gebrauchende Uhr um $13'\,48''$ vor Göttinger mittlerer Zeit voraus sei, so würden die Bestimmungen des Standes der Nadel für die Uhrzeiten $0^{\mathrm h}\,14'\,.\ .\ .\ .\ .\allowbreak\,0^{\mathrm h}\,19'\,.\ .\ .\ .\ .\allowbreak\,0^{\mathrm h}\,24'\,.\ .\ .\ .\ .\allowbreak\,0^{\mathrm h}\,29'$ u.\ s.\ f.\ zu machen sein. \textit{Volle} Minuten zu wählen, ist aber jedenfalls anzurathen, weil man sich so die einzelnen Operationen leichter mechanisch macht.

Unter dem Stand der Magnetnadel, welcher für die einzelnen Zeitmomente bestimmt werden soll, ist hier nicht diejenige Stellung verstanden, welche der aufgehängte Magnetstab in dem betreffenden Augenblick wirklich eben hat, sondern diejenige, welche er haben würde, wenn er (oder genauer zu reden, seine magnetische Axe) in diesem Augenblick genau im magnetischen Meridian wäre. Diese Distinction war unnöthig, so lange man sich nur solcher Nadeln bediente, die eine sehr grosse Genauigkeit nicht geben konnten: man brauchte nur dafür zu sorgen, dass die Nadel um die Zeit der Beobachtung in keiner erkennbaren Schwingung begriffen war, und erhielt damit das Gesuchte unmittelbar. Bei den viel grössern Forderungen, die man an die Genauigkeit der Bestimmungen durch die jetzt eingeführten Apparate machen kann und machen muss, kann aber von einer solchen unmittelbaren Bestimmung nicht mehr die Rede sein. Es steht nicht in unsrer Macht, die Nadel des Magnetometers \textit{so} vollkommen zu beruhigen, dass gar keine erkennbaren Schwingungsbewegungen zurückbleiben; wenigstens kann es nicht mit Sicherheit ohne Zeitaufwand, und nicht auf die Dauer geschehen. Es werden daher an die Stelle der unmittelbaren Beobachtung solche mittelbare Bestimmungen treten müssen, zu denen eine vollkommene Beruhigung unnöthig ist.

Die sich zuerst darbietende Methode besteht darin, dass man die Nadel absichtlich im schwingenden Zustande beobachtet, zwei auf einander folgende äusserste Stellungen (ein Minimum und ein Maximum) an der Scale aufzeichnet, und zwischen beiden das Mittel nimmt. Dieses an sich unverwerfliche Verfahren erfordert jedoch, wenn die Schwingungen eine beträchtliche Grösse haben, eine Modification, und ist, wenn die Schwingungen klein sind, nur unter einer einschränkenden Bedingung zulässig. Im ersten Fall nemlich wird selbst von einer Schwingung zur andern die successive Abnahme des Schwingungsbogens nicht unmerklich, daher auch schon die Abweichung vom wirklichen Meridian auf der Maximum-Seite geringer sein, als sie beim vorhergehenden Minimum auf der ent-
\par\clearpage
\SourceLeaf{549}\GaussVSourceFolio{543}{R}{MAGNETISCHE BEOBACHTUNGEN.}
\noindent gegengesetzten Seite gewesen war, folglich das Mittel aus diesem Minimum und dem folgenden Maximum zu klein werden. Aus derselben Ursache wird das Mittel aus diesem Maximum und dem folgenden Minimum ein zu grosses Resultat geben. Da nun aber die Abnahme des Schwingungsbogens einige Schwingungen hindurch beinahe gleichförmig bleibt, so kann man das Mittel aus zwei solchen Mitteln als hinlänglich genau, und zwar als geltend für den Augenblick der zweiten Elongation betrachten. Oder, um es durch eine Formel auszudrücken, wenn $a,b,c$ die Ablesungen in drei auf einander folgenden Elongationen sind (gleich viel, ob die erste und dritte Minima sind, und die zweite ein Maximum, oder umgekehrt), so stellt $\tfrac14(a+2b+c)$ den im Augenblick der Elongation $b$ Statt findenden Stand des magnetischen Meridians dar.

Bei kleinen Schwingungen ist dieses Verfahren nur dann zulässig, wenn die Declination keinen in kurzer Zeit merklichen Veränderungen unterworfen ist, und man kann dann schon das Mittel aus zwei auf einander folgenden Elongationen, als für den in der Mitte liegenden Augenblick gültig ansetzen: im entgegengesetzten Fall aber, d.\ i.\ zu einer Zeit, wo in der Declination schnell beträchtliche Änderungen vorgehen, kann dies Verfahren seine Brauchbarkeit gänzlich verlieren.

Immer aber behält die Methode, den Stand des magnetischen Meridians aus beobachteten Elongationen zu bestimmen, die Unbequemlichkeit, dass die Augenblicke, für welche das erhaltene Resultat gilt, nicht dieselben sind (oder es nur zufällig werden), für welche man den Stand verlangt. Und wenn auch dies in der Mehrzahl der Fälle wenig erheblich sein mag, so verdient doch offenbar ein anderes Verfahren den Vorzug, welches, von jener Inconvenienz frei, Bequemlichkeit, Gleichförmigkeit und alle nur zu wünschende Schärfe in sich vereinigt, und deshalb von sämmtlichen Theilnehmern an den Terminsbeobachtungen befolgt wird.

Dieses Verfahren beruht auf dem Satze, dass das Mittel aus zwei Stellungen der Nadel, die zweien genau um eine Schwingungsdauer von einander abstehenden Augenblicken entsprechen, mit derjenigen Lage des magnetischen Meridians übereinstimmt, welche für das Mittel dieser Zeiten Statt fand, in welche Theile der Schwingungsperiode diese Zeiten auch fallen mögen. Dieser Satz würde in mathematischer Schärfe wahr sein, wenn theils keine äussere Ursachen (wie der Widerstand der Luft u.\ dergl.) zur successiven Verkleinerung~des~Schwingungs-
\par\clearpage
\begingroup
\renewcommand{\footnoterule}{\kern-4.35pt\hrule width .256\columnwidth height .35pt\kern4pt}
\SourceLeaf{550}\GaussVSourceFolio{544}{L}{MAGNETISCHE BEOBACHTUNGEN.}
\noindent bogens. wirkten, theils die etwaige Veränderung in der Lage des magnetischen Meridians während jener kurzen Zwischenzeit nur als \textit{gleichförmig} betrachtet werden dürfte. Der erstere Umstand hat aber gar keinen merklichen Einfluss, wenn man das Verfahren immer nur auf sehr kleine Schwingungsbewegungen anwendet, und was den zweiten betrifft, so sind die Veränderungen der Declination während einer so kurzen Zwischenzeit in der Regel schon an sich kaum merklich, und um so mehr ist man berechtigt, wenigstens die Gleichförmigkeit der Veränderungen während dieser kurzen Zeit gelten zu lassen\HistoricalNote{550}.

Hiemit ist nun die Aufgabe von selbst gelöst. Um den der Declination für die Zeit $T$ entsprechenden Stand der Nadel zu erfahren, braucht man nur, nachdem nöthigenfalls vorher ihre Bewegungen durch angemessene Beruhigungsmittel auf sehr kleine gebracht sind, die wirklichen Stellungen für die Zeiten $T-\tfrac12t$ und $T+\tfrac12t$ zu beobachten, und daraus das Mittel zu nehmen, wo $t$ die Schwingungsdauer bedeutet. Inzwischen, grösserer Genauigkeit und Sicherheit wegen, beschränkt man sich hierauf nicht, sondern macht noch einige ähnliche Bestimmungen für ein Paar Zeitmomente kurz vor, und eben so viele nach $T$, immer in gleichen Intervallen, unter welcher Voraussetzung, insofern während dieser Zeit die Änderung der Declination als gleichförmig betrachtet werden darf, das Mittel aus allen diesen Resultaten das für die Zeit $T$ geltende \textit{Endresultat} sein wird, und zuverlässiger als die einzelne Bestimmung für $T$ selbst.

Die einfachste Art, dies auszuführen, besteht, wenn z.\ B.\ das Endresultat auf fünf partiellen Resultaten beruhen soll, darin, dass man den wirklichen Stand der Nadel für die sechs Zeiten
\[
T-\tfrac52t,\quad T-\tfrac32t,\quad T-\tfrac12t,\quad T+\tfrac12t,\quad T+\tfrac32t,\quad T+\tfrac52t
\]
aufzeichnet. Sind die aufgezeichneten Zahlen $a,b,c,d,e,f$, so wird $\tfrac12(a+b)$ das für die Zeit $T-2t$ geltende Resultat sein; eben so
\[
\tfrac12(b+c),\quad\tfrac12(c+d),\quad\tfrac12(d+e),\quad\tfrac12(e+f)
\]
für die Zeiten $T-t$, $T$, $T+t$, $T+2t$; und das Mittel aus diesen partiellen
\par\clearpage
\endgroup
\SourceLeaf{551}\GaussVSourceFolio{545}{R}{MAGNETISCHE BEOBACHTUNGEN.}
\noindent Resultaten oder der fünfte Theil ihrer Summe wird als berichtigtes Endresultat für die Zeit $T$ anzunehmen sein.

\Needspace{15\baselineskip}
Als Beispiel möge hier das Detail der Beobachtung in Göttingen am 17. August 1836 für $15^{\mathrm h}\,30'$ stehen. Der Beobachter war Hr.\ Dr.\ \textsc{Wappäus}. Für $t$ war angenommen $20''$.
\begingroup
\renewcommand{\arraystretch}{1.12}
\[
\begin{array}{@{}r@{}l@{\hspace{.40em}}r@{}l@{\hspace{.40em}}r@{}l|r|}
15&{}^{\mathrm h}&29&{}'&10&{}''&865.2\\
&&&&30&&867.5\\
&&&&50&&866.2\\
&&30&&10&&868.0\\
&&&&30&&867.3\\
&&&&50&&868.5
\end{array}
\quad
\left.\begin{array}{@{}r@{}}866.35\\866.85\\867.10\\867.65\\867.90\end{array}\right\}\,867.16
\]
\endgroup

Die erste Columne enthält hier die Beobachtungszeiten, die zweite die aufgezeichneten Scalentheile, die dritte das Mittel zwischen je zwei auf einander folgenden Aufzeichnungen, mithin die für
\[
15^{\mathrm h}\,29'\,20'',\quad15^{\mathrm h}\,29'\,40'',\quad15^{\mathrm h}\,30'\,0'',\quad15^{\mathrm h}\,30'\,20'',\quad15^{\mathrm h}\,30'\,40''
\]
geltenden partiellen Resultate, und daneben das für $15^{\mathrm h}\,30'\,0''$ geltende Endresultat. In diesem Beispiele ist die im Laufe der Beobachtungen fortwährend Statt habende Veränderung der Declination offenbar, und wird auch durch die vorhergehenden und folgenden Resultate bestätigt. Es war nemlich das Resultat
\[
\begin{array}{@{}lr@{}l@{\hspace{.40em}}r@{}l@{\hspace{.40em}}r@{}lll@{}}
\text{für}&15&{}^{\mathrm h}&25&{}'&0&{}''&.\ .\ .\ .\ .\ .&862.82\\
&&&35&&0&&.\ .\ .\ .\ .\ .&872.32
\end{array}
\]

Gewöhnlicher übrigens, als so beträchtliche Änderungen, ist der während der Zeit, welche ein Beobachtungssatz erfordert, fast stationäre Stand der Declination, und in solchen Fällen dient das kleinere oder grössere Hinundherschwanken der partiellen Resultate als ein Maassstab für die grössere oder geringere Zuverlässigkeit der Beobachtungen selbst, möge sie nun von dem Grad der Geschicklichkeit und Aufmerksamkeit des Beobachters, oder der Güte des Apparats selbst, oder von den mehr oder weniger günstigen äussern Umständen abhängen.

Das beschriebene Verfahren ist dasjenige, welches die meisten Theilnehmer an den Terminsbeobachtungen befolgen. Es setzt die Kenntniss der Schwingungsdauer der Nadel voraus, welche bekanntlich zugleich von~der~Stärke~der
\GaussVSourceGathering{69}
